% TEMPLATE-MILC-v3.tex - plantilla maestra para todas las guias MILC v3
% Ver scripts/generadores/build-guia-milc.py para la lista completa de
% claves esperadas. El script valida que no queden placeholders sin
% reemplazar antes de escribir el archivo final.

\documentclass[11pt,letterpaper]{article}
\usepackage[margin=1.75cm]{geometry}
\usepackage[table]{xcolor}
\usepackage{fontspec}
\usepackage[spanish]{babel}
\usepackage{tikz}
\usepackage[most]{tcolorbox}
\usepackage{tabularx}
\usepackage{array}
\usepackage{fancyhdr}
\usepackage{titlesec}
\usepackage{ragged2e}
\usepackage{enumitem}
\usepackage{microtype}

\setmainfont{Helvetica Neue}
\setsansfont{Helvetica Neue}
\setlength{\parindent}{0pt}
\setlength{\parskip}{6pt}
\hyphenpenalty=200
\exhyphenpenalty=200
\pretolerance=400
\tolerance=2000
\setlength{\emergencystretch}{1em}
\renewcommand{\arraystretch}{1.25}
\setlist{leftmargin=5mm,itemsep=1mm,topsep=1mm}
\newcolumntype{Y}{>{\RaggedRight\arraybackslash}X}

\definecolor{milcMagenta}{HTML}{E6007E}
\definecolor{milcVino}{HTML}{5A0038}
\definecolor{milcTurquesa}{HTML}{0093A5}
\definecolor{milcVerde}{HTML}{58A923}
\definecolor{milcMostaza}{HTML}{E5B400}
\definecolor{milcOcre}{HTML}{B86600}
\definecolor{milcNegro}{HTML}{241816}
\definecolor{milcGris}{HTML}{F7F7F4}
\definecolor{milcLinea}{HTML}{E8E2DD}
\definecolor{milcGradePrimary}{HTML}{84CC16}
\definecolor{milcGradeDark}{HTML}{4D7C0F}
\definecolor{milcGradeSoft}{HTML}{ECFCCB}
\definecolor{milcGradeAccent}{HTML}{CFE7FF}
\definecolor{milcGradeText}{HTML}{FFFFFF}
\color{milcNegro}

\pagestyle{fancy}
\fancyhf{}
\lhead{\small Tecnología e Informática · Grado 7}
\rhead{\small Guía 28 · MILC v3}
\cfoot{\small\thepage}
\renewcommand{\headrulewidth}{0pt}

\newtcolorbox{titlebox}[1]{
  enhanced,colback=#1,colframe=#1,arc=7mm,boxrule=0pt,
  left=7mm,right=7mm,top=3mm,bottom=3mm,
  before skip=8pt,after skip=8pt,
  fontupper=\sffamily\bfseries\Large\color{white}
}

\newtcolorbox{softbox}[3]{
  enhanced,colback=#2,colframe=#1,borderline west={4pt}{0pt}{#1},
  arc=2mm,boxrule=.4pt,left=4mm,right=4mm,top=3mm,bottom=3mm,
  before skip=6pt,after skip=8pt,title={#3},coltitle=white,
  colbacktitle=#1,fonttitle=\sffamily\bfseries,fontupper=\RaggedRight,
  attach boxed title to top left={xshift=3mm,yshift=-2mm},
  boxed title style={arc=2mm, boxrule=0pt}
}

\newtcolorbox{drawbox}[2]{
  enhanced,colback=white,colframe=#1,arc=4mm,boxrule=1pt,
  left=5mm,right=5mm,top=4mm,bottom=4mm,before skip=8pt,after skip=8pt,
  title={#2},coltitle=white,colbacktitle=#1,fonttitle=\sffamily\bfseries,
  fontupper=\RaggedRight,
  attach boxed title to top left={xshift=4mm,yshift=-2mm},
  boxed title style={arc=3mm, boxrule=0pt}
}

\newtcolorbox{infoband}[1]{
  enhanced,colback=#1!8,colframe=#1!50,arc=2mm,boxrule=.5pt,
  left=4mm,right=4mm,top=2mm,bottom=2mm,before skip=4pt,after skip=4pt,
  fontupper=\small\RaggedRight
}

% Puente narrativo entre secciones (contrato editorial Conéctate).
% Da continuidad: "Acabas de X. Ahora vas a Y porque Z."
% Visualmente: banda izquierda en color del grado, texto en itálica pequeña.
\newtcolorbox{puentebox}{
  enhanced,colback=milcGradeSoft,colframe=milcGradeDark,
  borderline west={3pt}{0pt}{milcGradeDark},
  arc=0mm,boxrule=0pt,left=5mm,right=4mm,top=2.5mm,bottom=2.5mm,
  before skip=4pt,after skip=8pt,
  fontupper=\itshape\small\color{milcNegro}\RaggedRight
}
\newcommand{\puente}[1]{%
  \begin{puentebox}%
    \textcolor{milcGradeDark}{\textbf{\textrightarrow}}\hspace{2mm}#1%
  \end{puentebox}%
}

\newcommand{\linead}[1]{\noindent\color{milcLinea}\rule{#1}{0.5pt}\color{milcNegro}}
\newcommand{\respuesta}[1]{\par\vspace{2mm}\foreach \n in {1,...,#1}{\noindent\color{milcLinea}\rule{\linewidth}{0.5pt}\par\vspace{5mm}}\color{milcNegro}}
\newcommand{\checkbox}{\textcolor{milcGradeDark}{\fbox{\phantom{x}}}\hspace{5pt}}

\newcommand{\phasecard}[4]{
\begin{tcolorbox}[enhanced,colback=#3,colframe=#2,arc=3mm,boxrule=.5pt,height=3.05cm,valign=center,halign=center,left=1.5mm,right=1.5mm,top=2mm,bottom=2mm]
{\sffamily\bfseries\fontsize{22}{24}\selectfont\textcolor{#2}{#1}}\par
{\sffamily\bfseries\small #4}
\end{tcolorbox}
}

\begin{document}
\thispagestyle{empty}
\begin{tikzpicture}[remember picture,overlay]
  \fill[white] (current page.south west) rectangle (current page.north east);
  \path[fill=milcGradePrimary,rounded corners=16mm]
    ([xshift=.55cm,yshift=-.55cm]current page.north west) rectangle
    ([xshift=-.55cm,yshift=3.65cm]current page.south east);

  \node[anchor=north west,text=milcGradeText] at ([xshift=2.0cm,yshift=-1.85cm]current page.north west)
    {\sffamily\bfseries\fontsize{14}{16}\selectfont GRADO};
  \node[anchor=north west,text=milcGradeText,opacity=.82] at ([xshift=2.0cm,yshift=-2.75cm]current page.north west)
    {\sffamily\bfseries\fontsize{9}{11}\selectfont SERIE GUÍAS · TECNOLOGÍA E INFORMÁTICA};

  \begin{scope}[shift={([xshift=-3.05cm,yshift=-2.55cm]current page.north east)}]
    \fill[white,opacity=.80] (-.62,-.52) rectangle (.62,.52);
    \draw[milcGradeText,line width=1pt] (-.62,-.52) rectangle (.62,.52);
    \fill[milcGradeDark] (-.42,-.38) rectangle (-.22,.06);
    \fill[milcGradeAccent] (-.10,-.38) rectangle (.10,.24);
    \fill[milcGradePrimary!60!black] (.22,-.38) rectangle (.42,.38);
  \end{scope}

  \node[anchor=west,text=milcGradeText] at ([xshift=1.9cm,yshift=-12.85cm]current page.north west)
    {\sffamily\bfseries\fontsize{124}{124}\selectfont 7};
  \draw[milcGradeText,line width=1.7pt]
    ([xshift=4.52cm,yshift=-11.68cm]current page.north west) circle (.13cm);

  \node[anchor=west,text=milcGradeText,text width=8.7cm,align=left] at ([xshift=10.35cm,yshift=-11.6cm]current page.north west)
    {
     {\sffamily\bfseries\fontsize{19}{22}\selectfont Ética de\\la IA ---\\sesgo,\\privacidad,\\derechos y\\deepfakes}\\[4mm]
     {\sffamily\bfseries\fontsize{23}{26}\selectfont Guía 28-7-TIC}\\[2mm]
     {\sffamily\fontsize{13}{16}\selectfont Período 3 · Inteligencia Artificial}\\[5mm]
     {\sffamily\bfseries\fontsize{11}{13}\selectfont Institución Educativa Sor María Juliana}\\[1mm]
     {\sffamily\fontsize{10}{12}\selectfont Autor: Dr. Álvaro Cárdenas Orozco}
    };

  \path[fill=milcGradeSoft,rounded corners=7mm]
    ([xshift=1.35cm,yshift=.85cm]current page.south west) rectangle
    ([xshift=-1.35cm,yshift=4.25cm]current page.south east);
  \draw[milcGradeDark,line width=.65pt,rounded corners=7mm]
    ([xshift=1.35cm,yshift=.85cm]current page.south west) rectangle
    ([xshift=-1.35cm,yshift=4.25cm]current page.south east);
  \node[anchor=center,text=milcNegro,text width=17.0cm,align=center] at ([yshift=2.55cm]current page.south)
    {\sffamily\fontsize{10}{12}\selectfont Metodología MILC · Apertura ancestral · Pensamiento computacional · Triángulo Dussel-Estoicismo-Floridi\\
    \textcolor{milcGradeDark}{\bfseries Producto:} En el cuaderno: \textbf{(a) 5 dilemas éticos analizados} (uno por cada tema: sesgo, privacidad, derechos de autor, deepfakes, autenticidad académica), con tu posición personal sobre cada uno y justificación; \textbf{(b) tu ``código personal de ética con IA''} con mínimo 5 reglas que prometes cumplir, firmado como compromiso; \textbf{(c) reflexión}: ``¿Qué tipo de ciudadano digital quiero ser cuando sea adulto?''.};
\end{tikzpicture}
\mbox{}
\newpage

\section*{Apertura: Ética de la IA --- sesgo, privacidad, derechos y deepfakes}

\begin{softbox}{milcGradeDark}{milcGradeSoft}{Saber ancestral}
\textbf{En el barrio de Cartago, la comunidad tenía reglas no escritas para usar herramientas poderosas.} Cuando alguien tenía un machete (herramienta para cortar caña, abrir camino, defenderse), \emph{no lo usaba para cualquier cosa}: lo respetaba. Si veías a un niño jugando con machete, le decías: \emph{``Eso no se juega; le puedes hacer daño a alguien''}. Si alguien usaba el machete para amenazar, la comunidad lo censuraba. \textbf{El machete no era malo en sí mismo}: lo malo era usarlo sin ética. \emph{Doña Mercedes, doña Esperanza la partera, los consejeros del barrio enseñaban estas reglas no escritas a los niños desde temprano}. La IA es la herramienta poderosa de tu generación. Como el machete, no es buena ni mala en sí misma: \textbf{depende de cómo la uses}. Hoy aprendes 5 dilemas éticos que ya están aquí, no en el futuro lejano. Y construyes tu propio código personal de ética con IA.
\end{softbox}

\begin{softbox}{milcTurquesa}{milcGris}{Saber contemporáneo}
La \textbf{ética de la IA} estudia los problemas morales que aparecen al usar IA. Hay \textbf{5 grandes temas} que afectan tu vida diaria. \textbf{(1) Sesgo (ya viste en S4):} la IA puede discriminar grupos por datos parciales. \textbf{(2) Privacidad:} cuando le das información a una IA, ¿dónde queda? ¿quién la lee? ¿la usan para entrenarse? \textbf{(3) Derechos de autor:} ChatGPT escribió un cuento. ¿De quién es? ¿Del usuario? ¿De OpenAI? ¿De los autores cuyos textos se usaron para entrenar la IA? \textbf{(4) Deepfakes:} videos o audios falsos generados por IA donde una persona dice o hace algo que NUNCA hizo. Pueden ser usados para acosar, estafar, manipular elecciones. \textbf{(5) Autenticidad académica:} usar IA para hacer tareas. ¿Es trampa? ¿Es asistencia legítima? Depende. \textbf{No hay respuestas fáciles}: estos son debates abiertos en el mundo. Tu generación va a decidir cómo se resuelven en los próximos 20 años. Hoy empiezas a tener opinión propia.
\end{softbox}

\begin{softbox}{milcMostaza}{milcGris}{Pregunta-puente}
\textit{Si tu amigo te muestra un \emph{deepfake} (video falso) de un político del barrio diciendo algo que nunca dijo, y te dice \emph{``mira lo que dijo''}, ¿\textbf{qué haces}? ¿Lo crees? ¿Lo compartes? ¿Cómo te das cuenta de que es falso?}
\end{softbox}

\begin{softbox}{milcVerde}{milcGris}{Saber y saber hacer}
Al terminar la clase: (1) podrás \textbf{identificar} los 5 temas éticos clave; (2) sabrás \textbf{analizar} dilemas éticos reales; (3) podrás \textbf{evaluar} situaciones con criterio propio; (4) habrás \textbf{creado} tu código personal de ética con IA.
\end{softbox}

\small
\begin{tabularx}{\textwidth}{>{\bfseries}p{3.0cm}Y}
\rowcolor{milcGradePrimary!12}Autor & Dr. Álvaro Cárdenas Orozco\\
DBA & Reconoce los fundamentos, usos y límites éticos de la Inteligencia Artificial (MEN, T\&I 7°)\\
Referentes & Saber ancestral del consejero · Modelos de lenguaje · Ética algorítmica y prompting\\
Producto final & En el cuaderno: \textbf{(a) 5 dilemas éticos analizados} (uno por cada tema: sesgo, privacidad, derechos de autor, deepfakes, autenticidad académica), con tu posición personal sobre cada uno y justificación; \textbf{(b) tu ``código personal de ética con IA''} con mínimo 5 reglas que prometes cumplir, firmado como compromiso; \textbf{(c) reflexión}: ``¿Qué tipo de ciudadano digital quiero ser cuando sea adulto?''.\\
\end{tabularx}
\normalsize

\puente{Hoy haces 4 cosas: identificas los 5 temas, analizas 5 dilemas reales, formulas tu posición, escribes tu código personal.}

\begin{titlebox}{milcGradeDark}
Ruta MILC: así vamos a aprender
\end{titlebox}
\begin{tabularx}{\textwidth}{>{\centering\arraybackslash}X>{\centering\arraybackslash}X>{\centering\arraybackslash}X>{\centering\arraybackslash}X}
\phasecard{1}{milcVerde}{milcGris}{Escuta\\[-1mm]\small Observo, escucho y pregunto.} &
\phasecard{2}{milcTurquesa}{milcGris}{Sistematización\\[-1mm]\small Pensamiento computacional.} &
\phasecard{3}{milcMagenta}{milcGris}{Praxis\\[-1mm]\small Construyo y aplico.} &
\phasecard{4}{milcVino}{milcGris}{Evaluación\\[-1mm]\small Reflexión liberadora.}
\end{tabularx}


\puente{Empezamos con 5 dilemas que ya existen, no son ciencia ficción.}

\begin{titlebox}{milcVerde}
Fase 1 · Escuta
\end{titlebox}

\begin{softbox}{milcVerde}{milcGris}{Escena para escuchar}
\textbf{Actividad 1 · EVALÚA --- 5 dilemas éticos reales} (12 min · individual). Lee estos 5 dilemas y en el cuaderno escribe tu posición + justificación (3-4 líneas por cada uno). \emph{Dilema 1 (sesgo):} \emph{una IA aprueba créditos bancarios y rechaza más a personas de barrios populares de Cartago. ¿Es justo? ¿Qué debería pasar?}. \emph{Dilema 2 (privacidad):} \emph{le cuentas a ChatGPT que tienes una pelea con tu papá y le pides consejo. ¿Está bien? ¿Quién más podría ver esa conversación?}. \emph{Dilema 3 (derechos de autor):} \emph{usas DALL-E para generar una imagen y la imprimes para venderla. ¿Es tuya la imagen? ¿De quién? ¿Es ético venderla?}. \emph{Dilema 4 (deepfake):} \emph{ves un video viral del presidente diciendo algo escandaloso. Resulta ser deepfake. ¿Qué consecuencias tiene esto para la democracia? ¿Cómo me protejo?}. \emph{Dilema 5 (autenticidad académica):} \emph{tu compañero usa ChatGPT para escribir su ensayo completo y lo entrega como suyo. ¿Es trampa? ¿O es como usar Wikipedia?}.
\end{softbox}

\checkbox Leí los 5 dilemas con calma.\\
\checkbox Pensé en mi posición personal.\\
\checkbox Escribí 3-4 líneas por cada uno con justificación.

\begin{infoband}{milcVerde}
\textbf{Lo que va al cuaderno (Actividad 1 · EVALÚA).} Título: «Actividad 1 --- 5 dilemas éticos». \textbf{Formato:} 5 bloques (uno por dilema), cada uno con mi posición + justificación. \textbf{Extensión:} 15-20 líneas. \textbf{Sabes que terminaste cuando} tienes opinión propia en cada uno, no copiada.
\end{infoband}

\puente{Después aprendes los 5 temas a fondo con casos reales.}

\begin{titlebox}{milcTurquesa}
Fase 2 · Sistematización con pensamiento computacional
\end{titlebox}

\begin{softbox}{milcTurquesa}{milcGris}{Cuatro pilares aplicados al tema}
Tu posición personal es válida. Ahora profundiza en cada uno de los 5 temas con casos reales documentados.
\end{softbox}

\small
\begin{tabularx}{\textwidth}{>{\bfseries\RaggedRight\arraybackslash}p{3.6cm}Y}
\rowcolor{milcTurquesa!90}\color{white}Pilar & \color{white}\bfseries Aplicación al tema\\
1. Descomposición & \textbf{Tema 1 --- SESGO ALGORÍTMICO.} Ya lo viste en S4. Resumen: las IAs reproducen los sesgos de sus datos. Casos reales: Amazon descartando CVs de mujeres, Google clasificando personas afro como animales, COMPAS sesgada en el sistema de justicia EEUU. \textbf{Tu responsabilidad ética como usuario:} (a) saber que existe; (b) verificar; (c) no aceptar pasivamente decisiones que parecen injustas; (d) exigir que los servicios que usas sean transparentes.\\
2. Reconocimiento de patrones & \textbf{Tema 2 --- PRIVACIDAD.} Cuando le escribes algo a ChatGPT, esa conversación \emph{queda guardada} en servidores de OpenAI. Algunas empresas (no todas) usan tus conversaciones para entrenar futuras versiones. Otras venden datos a terceros. \emph{No compartas info sensible:} cédula, contraseñas, dirección exacta, info médica, asuntos íntimos. \textbf{Una vez escrito, no sabes dónde queda}. Caso real: empleados de Samsung pegaron código secreto en ChatGPT (2023), el código quedó en los servidores de OpenAI. Samsung prohibió ChatGPT a sus empleados.\\
3. Abstracción & \textbf{Tema 3 --- DERECHOS DE AUTOR.} La IA aprendió de millones de textos e imágenes (muchas con derechos de autor). Cuando ChatGPT escribe un cuento, ¿es plagio? ¿Es creación nueva? \emph{Los tribunales del mundo están debatiendo esto}. \textbf{Tu posición práctica:} (a) atribuye cuando uses ideas de IA en tus trabajos; (b) no vendas creaciones de IA como tuyas sin atribución; (c) reconoce que la cosa creativa ``híbrida'' (parcialmente tuya, parcialmente IA) es categoría nueva.\\
4. Algoritmo conceptual & \textbf{Tema 4 --- DEEPFAKES.} Videos, imágenes y audios falsos generados por IA. Tan convincentes que pueden engañar al ojo humano. \emph{Usos malos:} pornografía no consentida, suplantación de identidad, fraude, manipulación política. \emph{Usos legítimos:} efectos en cine, doblaje multilingüe, accesibilidad. \textbf{Cómo protegerse:} verificar fuente, buscar artefactos en bordes/manos, contrastar con fuentes confiables, NO compartir sin verificar. \textbf{Tema 5 --- AUTENTICIDAD ACADÉMICA.} ¿Es trampa usar IA en tareas? \emph{Depende del uso:} (a) usar IA para entender = LEGÍTIMO; (b) usar IA para estructurar = LEGÍTIMO con atribución; (c) usar IA para escribir el ensayo completo y entregarlo como tuyo = TRAMPA. La diferencia es \emph{si tú aprendiste o solo entregaste un producto}. Los colegios están definiendo políticas; los profesores también. \emph{La regla universal: si no aprendiste, es trampa}.\\
\end{tabularx}
\normalsize

\begin{drawbox}{milcTurquesa}{5 temas éticos clave}
\textbf{1.} SESGO: discriminación por datos parciales. \\[1mm]
\textbf{2.} PRIVACIDAD: tus datos quedan en servidores. \\[1mm]
\textbf{3.} DERECHOS DE AUTOR: la creación con IA es categoría nueva. \\[1mm]
\textbf{4.} DEEPFAKES: videos/audios falsos convincentes. \\[1mm]
\textbf{5.} AUTENTICIDAD ACADÉMICA: usar IA en tareas; depende del uso.
\end{drawbox}

\begin{softbox}{milcMostaza}{milcGris}{Errores comunes y su corrección}
\textbf{Errores éticos frecuentes con IA.} (1) \emph{Aceptar respuestas de IA sin verificar.} Los sesgos pasan a tu pensamiento sin que lo notes. (2) \emph{Compartir datos sensibles con IA.} Una vez escritos, no sabes dónde quedan. (3) \emph{Copiar de IA sin atribución.} Es la nueva forma de plagio. (4) \emph{Compartir deepfakes sin verificar.} Te haces cómplice de desinformación. (5) \emph{Usar IA para hacer la tarea sin aprender.} Te haces dependiente y pierdes el oficio.
\end{softbox}

\begin{infoband}{milcTurquesa}
\textbf{Lo que va al cuaderno (Actividad 2 · EXPLICA).} Título: «Actividad 2 --- 5 temas éticos + casos reales». \textbf{Formato:} tabla de 5 filas y 3 columnas (tema, caso real, mi responsabilidad). \textbf{Extensión:} 5 filas. \textbf{Sabes que terminaste cuando} puedes hablar de cada tema con criterio, no con opinión vacía.
\end{infoband}

\puente{Con todo claro, escribes tu código personal de ética + análisis de dilemas.}

\begin{titlebox}{milcMagenta}
Fase 3 · Praxis con pensamiento computacional
\end{titlebox}

\begin{softbox}{milcMagenta}{milcGris}{Construyendo el producto con los cuatro pilares}
\textbf{Actividad 3 · CREA --- Mi código personal de ética con IA} (28 min · individual). Vas a construir tu propio código de ética: 5+ reglas que prometes cumplir. Es ejercicio adulto de definir cómo vas a usar IA durante tu vida.
\end{softbox}

\small
\begin{tabularx}{\textwidth}{>{\bfseries\RaggedRight\arraybackslash}p{3.6cm}Y}
\rowcolor{milcMagenta!90}\color{white}Pilar & \color{white}\bfseries Aplicación al producto\\
Descomposición & \textbf{Paso 1 --- Título y dedicatoria.} En el cuaderno escribe: \textbf{``Mi código personal de ética con IA''}. Debajo, en una línea: \emph{``Comprometido con mi yo de 25 años, [tu nombre], grado 7°, año 2026''}. Este código no es ``para entregar al profe''; es para ti.\\
Patrones & \textbf{Paso 2 --- Mis 5 reglas (mínimo).} Lista numerada del 1 al 5+. Sugerencias para inspirarte (no copies; adáptalas a tu voz): \emph{(1) Siempre verifico lo que la IA me dice antes de usarlo. (2) No comparto datos personales sensibles con ningún chatbot. (3) Atribuyo cuando uso IA como asistente en tareas: la cito como fuente. (4) Solo creo contenido con IA si lo identifico como tal (no lo paso como creación 100\% propia). (5) No comparto deepfakes ni contenidos falsos sin verificar primero. (6) Uso la IA para aprender, no para evitar aprender. (7) Reviso una vez al mes mis hábitos con IA: ¿estoy siendo coherente con mi código?}.\\
Abstracción & \textbf{Paso 3 --- Justificación de cada regla.} Debajo de cada una, escribe 1 línea explicando \emph{por qué la elegiste}. Esto te ayuda a memorizarla y le da sentido. Ejemplo: \emph{``Regla 1 elegida porque sé que las IAs alucinan y no quiero citar mentiras como verdades''}.\\
Algoritmo de armado & \textbf{Paso 4 --- Análisis profundo de los 5 dilemas de la Actividad 1.} Vuelve a tus 5 respuestas de la Actividad 1. Con lo aprendido en la sistematización, \emph{¿cambiarías alguna respuesta?}. Anota: \emph{``Después de aprender la sistematización, mi posición sobre el dilema [N] cambió porque \_\_\_''}. Si no cambia, está bien; lo importante es revisar con criterio.\\
\end{tabularx}
\normalsize

\begin{drawbox}{milcMagenta}{Antes de cerrar la S8}
\checkbox Mi código personal tiene mínimo 5 reglas. \\[1mm]
\checkbox Cada regla tiene 1 línea de justificación. \\[1mm]
\checkbox Revisé mis 5 dilemas de la Actividad 1 con la sistematización. \\[1mm]
\checkbox Hay reflexión de qué tipo de ciudadano digital quiero ser. \\[1mm]
\checkbox El código tiene firma de cierre comprometida. \\[1mm]
\checkbox Lo guardo como documento personal para revisar anualmente.
\end{drawbox}

\begin{softbox}{milcMagenta}{milcGris}{Plantilla de trabajo}
\textbf{Estructura.} \textit{Paso 1:} título + dedicatoria a mi yo de 25. \textit{Paso 2:} mínimo 5 reglas. \textit{Paso 3:} justificación de cada una. \textit{Paso 4:} revisión de dilemas con sistematización. \textit{Paso 5:} reflexión ciudadano digital. \textit{Paso 6:} firma comprometida.
\end{softbox}

\begin{infoband}{milcMagenta}
\textbf{Lo que va al cuaderno (Actividad 3 · CREA).} Título: «Actividad 3 --- Mi código personal de ética con IA». \textbf{Formato:} 5+ reglas + justificación + revisión + reflexión + firma. \textbf{Extensión:} 1 página completa. \textbf{Sabes que terminaste cuando} sientes que el código refleja tu yo verdadero y lo cumplirás.
\end{infoband}

\puente{El producto es ese código firmado + 5 dilemas analizados.}

\begin{titlebox}{milcMostaza}
Producto de la sesión
\end{titlebox}
\textbf{Mi código personal de ética con IA · 5+ reglas + reflexión · firmado}.
\begin{softbox}{milcMostaza}{milcGris}{Criterios de calidad}
(1) El código tiene mínimo 5 reglas personales (no copiadas). (2) Cada regla tiene justificación de por qué la elegiste. (3) Revisaste los 5 dilemas de la Actividad 1 con criterio nuevo. (4) Hay reflexión sobre qué tipo de ciudadano digital quieres ser. (5) El código está firmado como compromiso real.
\end{softbox}

\puente{Antes de salir, comparte tu código con un compañero --- ¿se parecen o son distintos?.}

\begin{titlebox}{milcVino}
Fase 4 · Evaluación liberadora — cinco dimensiones
\end{titlebox}

\begin{softbox}{milcVino}{milcGris}{Sentido de la evaluación}
Evaluar no es solo revisar una nota. Es reconocer qué aprendí, qué cambió en mí, qué emociones manejé, cómo me relaciono con la ciudad y la comunidad, y qué le dejaré a quien viene detrás.
\end{softbox}

\textbf{Mi reflexión liberadora en cinco dimensiones.} Completa cada frase iniciada:

\small
\begin{tabularx}{\textwidth}{>{\bfseries\RaggedRight\arraybackslash}p{3.4cm}Y>{\RaggedRight\arraybackslash}p{4.6cm}}
\rowcolor{milcVino}\color{white}Dimensión & \color{white}\bfseries Mi respuesta (frame starter) & \color{white}\bfseries Algo que descubrí\\
1. Personal & Lo que más cambió en mí fue \linead{6.5cm} & \linead{4.2cm}\\
2. Emocional & La emoción más fuerte fue \linead{2.5cm} y la regulé \linead{3cm} & \linead{4.2cm}\\
3. Ciudadana & En democracia esto importa porque \linead{6cm} & \linead{4.2cm}\\
4. Local (Cartago) & En mi barrio sería útil para \linead{6.5cm} & \linead{4.2cm}\\
5. Intergeneracional & A mi abuela/abuelo le diría \linead{6.5cm} & \linead{4.2cm}\\
\end{tabularx}
\normalsize

\begin{infoband}{milcVino}
\textbf{Para la próxima sustentación.} Vuelve a esta tabla en seis meses. Verás que las respuestas cambian — no porque lo aprendido cambie, sino porque tú cambias. La evaluación liberadora es una conversación contigo a través del tiempo.
\end{infoband}


\puente{Cerramos con 3 voces sobre por qué la ética con IA es responsabilidad de tu generación.}

\begin{titlebox}{milcGradeDark}
Triángulo de pensamiento · cierre filosófico
\end{titlebox}

\begin{softbox}{milcVino}{milcGris}{Dussel · lente del nosotros}
\textit{````La ética no es lujo de unos pocos: es responsabilidad de todos. Cuando usas IA, tus decisiones pequeñas suman a un mundo más justo o menos.''''}

\textbf{Cada vez que decides usar IA con criterio} --- verificar, no compartir datos, atribuir, no propagar deepfakes --- estás aportando a un mundo digital más justo. \emph{Cada vez que decides usarla sin ética}, estás aportando a uno menos justo. No es exageración: las decisiones individuales se suman. \emph{Tu código personal hoy es voto por la sociedad digital que quieres habitar en 20 años}.

\textbf{¿Mi código refleja la sociedad digital en la que quiero vivir cuando sea adulto?}
\end{softbox}
\linead{\linewidth}\\[1mm]
\linead{\linewidth}

\begin{softbox}{milcMostaza}{milcGris}{Estoicismo · Marco Aurelio (emperador de la integridad personal) · cuidado interior}
\textit{````Sé el mismo cuando nadie te ve que cuando todos te miran. La integridad no se ajusta a la audiencia.''''}

\textbf{La integridad con IA es lo mismo}. Cuando le pides al LLM ayudar con tu tarea, \emph{nadie te observa}: solo tú y el chat. \emph{La decisión de usarlo con ética o sin ética la tomas tú solo}. El estoico hace lo correcto sin importar quién mira. Tu código personal te entrena en esa integridad antigua, aplicada al mundo digital.

\textbf{Cuando uso IA solo, ¿soy el mismo que cuando alguien me observa?}
\end{softbox}
\linead{\linewidth}\\[1mm]
\linead{\linewidth}

\begin{softbox}{milcTurquesa}{milcGris}{Floridi · lente de la infoesfera}
\textit{````En 2030 habrá millones de personas usando IA sin ética alguna. Los pocos que la usen con ética serán los líderes de la próxima década.''''}

\textbf{La mayoría de adultos hoy usa IA sin pensar en ética}: copia, plagia, propaga deepfakes, no verifica. Eso \emph{abre espacio para que los que sí piensan en ética destaquen}. \emph{En universidad, en trabajo, en proyectos sociales}: los que usan IA con criterio ético van a ser reconocidos. Tu código personal hoy es preparación para ese liderazgo futuro.

\textbf{¿Estoy entrenando una ética con IA que en 5 años pocos tendrán?}
\end{softbox}
\linead{\linewidth}\\[1mm]
\linead{\linewidth}

\begin{titlebox}{milcVino}
Compromiso final
\end{titlebox}

\small
\begin{tabularx}{\textwidth}{>{\RaggedRight\arraybackslash}p{3.0cm}YYY}
\rowcolor{milcVino}\color{white}\bfseries Criterio & \color{white}\bfseries Lo logré & \color{white}\bfseries En proceso & \color{white}\bfseries Necesito apoyo\\
Comprensión del tema & Explico ideas centrales con mis palabras. & Reconozco ideas pero falta precisión. & Necesito repasar conceptos base.\\
Pensamiento computacional & Apliqué los 4 pilares al diseñar mi producto. & Apliqué algunos pilares. & Necesito releer la fase 2.\\
Aplicación y producto & Mi producto cumple los criterios y se puede revisar. & Producto funcional con ajustes pendientes. & Aún en construcción.\\
Reflexión 5 dimensiones & Respondí cada dimensión con honestidad. & Respondí algunas con detalle. & Mis respuestas son muy generales.\\
Triángulo de pensamiento & Conecté las 3 voces con mi trabajo real. & Conecté al menos una. & No relaciono las voces con mi proceso.\\
\end{tabularx}
\normalsize

\textbf{Hoy entendí que:}\\[1mm]
\linead{\linewidth}\\[1mm]
\linead{\linewidth}

\textbf{Mi compromiso para la próxima sesión es:}\\[1mm]
\linead{\linewidth}\\[1mm]
\linead{\linewidth}

\begin{softbox}{milcMostaza}{milcGris}{Cita de despedida · Marco Aurelio}
\textit{``El obstáculo en el camino se vuelve el camino. Lo que estorba al camino se vuelve camino.''}\\[1mm]
Cada error de hoy, bien observado, se convierte en pieza del camino que pisarás mañana.
\end{softbox}

\small
\begin{tabularx}{\textwidth}{>{\bfseries}p{3.6cm}Y}
\rowcolor{milcGradeDark!12}Nombre completo: & \linead{10cm}\\
Fecha: & \linead{4cm} \quad Curso/grupo: \linead{4cm}\\
Mi firma: & \linead{9cm}\\
\end{tabularx}
\normalsize

\begin{infoband}{milcGradeDark}
\textbf{Para la próxima sesión.} Esta semana, aplica al menos 2 reglas de tu código personal en situaciones reales. Anota qué pasó. La próxima clase aplicamos todo a tu vida académica: \textbf{IA y mis tareas}. Vas a definir reglas concretas para usar IA en estudios.
\end{infoband}

\end{document}
